\clearpage
\chapter*{강연 221의 정오표}
\addcontentsline{toc}{chapter}{강연 221의 정오표}
\setcounter{footnote}{3}
\src{302}

\noindent\textbf{6. TDTE IV: 힐베르트 스킴.}

\begin{center}
[부르바키 세미나, 제13권, 1960/61, 제221호, 28쪽.]
\end{center}

\noindent\textbf{6쪽, 정리 2.1, 3행.} --- ``필요하다'' 대신
``필요충분하다''로 읽어야 한다.

\noindent\textbf{15쪽, 13행.} --- ``$P\neq Q$'' 대신
``$P<Q$''로 읽어야 한다.

\noindent\textbf{15행.} --- ``반대로 $P=Q$이면'' 대신
``그렇지 않은 경우에는 충분히 큰 $n$에 대하여 $P(n)\geq Q(n)$이다''로
읽어야 한다.

\noindent\textbf{18행, 공식 (*).} --- ``$<$''와 ``$>$'' 대신
``$\leq$''와 ``$\geq$''로 읽어야 한다.

\noindent\textbf{아래에서 5행.} --- ``$r$'' 대신 ``$Z$''로
읽어야 한다.

\noindent\textbf{18쪽, 주 3.9.} --- 대수적으로 닫힌 체 위의 힐베르트
스킴들의 연결 성분에 대한 연구는 하츠혼이 수행했음을 지적해 둔다. 그는
$\operatorname{Hilb}^{P}_{r}$가 연결임을 증명하고,
$\operatorname{Hilb}^{P}_{r}\neq\varnothing$이 되는 쌍 $(r,P)$를
결정한다\footnote{HARTSHORNE (R.). --- Connectedness of the Hilbert
scheme (Thesis, Harvard, 1962).}.

\noindent\textbf{23쪽, 아래에서 5행.} --- ``한 성분'' 대신
``한 성분의''로 읽어야 한다.

\noindent\textbf{23쪽과 24쪽, 주 5.5.} --- 자파의 예와 관련하여,
멈퍼드가 최근 $\mathbf P^3_{\mathbf C}$에 대한 힐베르트 스킴의 한 기약
성분을 구성하기까지 했음을 지적해 둔다. 이 성분의 일반적인 점들은 차수
14, 종수 24인 비특이 곡선들을 나타내며, 이 성분은 그 일반점에서
비축약이다. 이렇게 얻은 곡선들 가운데 하나를 중심으로
$\mathbf P^3_{\mathbf C}$를 블로업함으로써, 그는 또한 $\mathbf C$ 위의
3차원 정칙 사영 스킴을 얻는다. 그 형식적 모듈라이 스킴은
그 일반점에서 비축약이다. 동치인 표현으로, $\mathbf C$ 위의
복소해석기하학의 의미에서 그 국소 모듈라이 다양체(cf. 카르탕 세미나,
제13권, 1960/61)는 모든 점에서 비축약이다.
